LLVManim is implemented as a three-layer pipeline---\textbf{Ingestion}, \textbf{Transformation}, and \textbf{Render}---orchestrated by a dedicated \textbf{CLI} entry point. The high-level structure is shown in Fig.~\ref{fig:architecture_detail}. Layer boundaries are enforced at build time by \texttt{import-linter}, ensuring that no module imports from a downstream layer. The entire codebase is written in Python~3.12 and managed with \texttt{uv}.

% \begin{figure*}[htbp]
% \centerline{\includegraphics[height=0.6\textheight, keepaspectratio]{architecture_overview}}
% \caption{High-level three-layer architecture of LLVManim.}
% \label{fig:architecture_overview}
% \end{figure*}

\begin{figure*}[htbp]
\centerline{\includegraphics[height=0.88\textheight]{01_architecture_overview_vertical}}
\caption{Detailed module-level architecture. The Ingestion layer collects LLVM IR and metadata, the Transformation layer maps it to an intermediate scene graph, and the Rendering layer renders animations or exports artifacts.}
\label{fig:architecture_detail}
\end{figure*}

\subsubsection{CLI Entry Point}
The \texttt{llvmanim.cli} module is the user-facing entry point. It accepts LLVM~IR files and orchestrates the three pipeline layers via \texttt{argparse}:
\begin{minted}[
  autogobble,
  frame=single,
  bgcolor=listingbg,
  breaklines=true,
  fontsize=\footnotesize,
]{bash}
# Export JSON scene graph + Graphviz DOT/PNG
uv run llvmanim input.ll --json --draw --outdir out/
# Render basic stack animation
uv run llvmanim input.ll --animate --ir-mode basic --outdir out/
# Render rich IR spotlight with SSA panel
uv run llvmanim input.ll --animate --ir-mode rich-ssa --speed 1.5 --outdir out/
# CFG traversal animation
uv run llvmanim input.ll --cfg-animate --dot-cfg main.dot --import-trace trace.json --outdir out/
\end{minted}
\captionof{listing}{Actual CLI usage}
Fig.~\ref{fig:cli_dispatch} traces the complete CLI dispatch flow, showing how each command flag routes through the shared Ingest pipeline to the appropriate Transform and Render stages. When the tool encounters unsupported IR constructs, missing debug metadata, or malformed input, it emits a descriptive error message and continues processing the remainder of the program where possible.

\begin{figure*}[htbp]
\centerline{\includegraphics[width=\textwidth]{05_cli_dispatch}}
\caption{CLI dispatch flow. Each command flag (\texttt{-{}-json}, \texttt{-{}-draw}, \texttt{-{}-animate}, \texttt{-{}-cfg-animate}) shares the Ingest pipeline and branches to the appropriate Transform and Render stages.}
\label{fig:cli_dispatch}
\end{figure*}

\subsubsection{Ingestion Layer}
The Ingestion layer extracts program structure from LLVM IR. It uses the \texttt{llvmlite}\cite{llvmlite} Python binding to parse \texttt{.ll} files and also supports importing pre-computed data from JSON.

\begin{itemize}
    \item \textbf{\texttt{llvm\_events.py}:} Parses an LLVM module and emits a~\texttt{ProgramEventStream}---a typed sequence of~\texttt{IREvent} records. Each event carries the function name, block name, opcode, full instruction text, a~classified \texttt{EventKind} (one of~\texttt{alloca}, \texttt{load}, \texttt{store}, \texttt{binop}, \texttt{compare}, \texttt{call}, \texttt{ret}, \texttt{br}, or~\texttt{other}), its index within the function, debug line number, and operand list. Typed CFG edges are extracted directly from \texttt{llvmlite} terminator operands (\texttt{br}, \texttt{switch}, \texttt{invoke}, \texttt{indirectbr}, \texttt{callbr}) with true/false labels on conditional branches.

    \item \textbf{\texttt{cfg\_edge\_io.py}:} Serializes and deserializes \texttt{CFGEdge} lists to~JSON, enabling round-trip import/export of control-flow edges via dedicated CLI flags.

    \item \textbf{\texttt{analysis\_metadata\_io.py}:} Loads and saves dominator-tree and loop-analysis metadata as~JSON, supporting round-trip import and export via CLI flags.

    \item \textbf{\texttt{trace\_io.py}:} Reads and writes runtime path traces as JSON (\texttt{-{}-import-trace} / \texttt{-{}-export-trace}), used for overlay visualization in CFG traversal animations.

    \item \textbf{\texttt{dot\_layout.py}:} Parses a Graphviz DOT file produced by \texttt{opt -passes=dot-cfg} and computes a \texttt{DotLayout} with node positions and edge routes for the CFG animation scene.

    \item \textbf{\texttt{display\_lines.py}:} Cleans raw IR text into per-function line lists used by the rich animation panels.
\end{itemize}

\subsubsection{Transformation Layer}
The Transformation layer converts the ingest output into an intermediate scene graph and execution traces suitable for rendering. The complete data model spanning all three layers is shown in Fig.~\ref{fig:data_model}.

\begin{figure*}[!htbp]
\centerline{\includegraphics[width=1.15\textwidth]{04_data_model}}
\caption{Unified data model. Ingest types (\texttt{ProgramEventStream}, \texttt{IREvent}, \texttt{CFGEdge}, \texttt{DotLayout}) feed into Transform types (\texttt{SceneGraph}, \texttt{SceneNode}, \texttt{SceneEdge}, \texttt{AnimationCommand}, \texttt{TraceOverlay}) which are consumed by Render types (\texttt{CommandDrivenScene}, \texttt{StackRenderer}, \texttt{CFGRenderer}).}
\label{fig:data_model}
\end{figure*}

\begin{itemize}
    \item \textbf{\texttt{models.py}:} Defines the core data model---\texttt{SceneGraph}, \texttt{SceneNode}, \texttt{SceneEdge}, \texttt{CFGBlock}, \texttt{TraceOverlay}, and \texttt{BlockMetadata}. Also defines \texttt{AnimationCommand} and \texttt{ActionKind}, the typed animation primitives consumed by the Render layer. Nodes carry a~\texttt{kind} string and an open \texttt{properties} dict, with an optional \texttt{animation\_hint} string. \texttt{SceneGraph} embeds an ordered \texttt{commands} list alongside its nodes and edges.

    \item \textbf{\texttt{scene.py}:} The \texttt{build\_scene\_graph} function groups events by function and block, attaches typed CFG edges, assigns block roles, applies optional domtree/loop metadata, and populates \texttt{SceneGraph.commands} via the overlay builder. \texttt{build\_stack\_scene\_graph} produces a call-tree-ordered \texttt{SceneGraph} with a fully populated \texttt{commands} list, translating the parsed \texttt{\allowbreak ProgramEventStream} directly into typed \texttt{\allowbreak AnimationCommand} records (carrying \texttt{ActionKind}, \texttt{target}, and \texttt{params}) ready for consumption by the Render layer.

    \item \textbf{\texttt{trace.py}:} The \texttt{build\_execution\_trace} function produces a~list of \texttt{TraceStep} or~\texttt{RichTraceStep} records capturing the execution order with optional SSA result values for \texttt{binop}, \texttt{compare}, and~\texttt{load} instructions.
\end{itemize}

\subsubsection{Render Layer}
The Render layer renders the scene graph using Manim Community Edition\cite{manimce} and provides export functionality. All scenes extend \texttt{CommandDrivenScene}, a base class that iterates over the \texttt{SceneGraph.commands} list and dispatches each \texttt{AnimationCommand} to a registered handler. The complete type hierarchy---from Ingest types through Transform intermediates to the Render dispatch classes---is shown in Fig.~\ref{fig:data_model}.

\begin{itemize}
    \item \textbf{Animation Modes:} The tool supports three stack-rendering modes, all driven by \texttt{StackRenderer}---a \texttt{CommandDrivenScene} subclass. Each mode consumes the same \texttt{AnimationCommand} records (detailed in Fig.~\ref{fig:data_model}) but presents progressively more detail.
    \begin{itemize}
        \item \texttt{-{}-ir-mode basic}: Stack-only view with per-cell badge flash for each instruction.
        \item \texttt{-{}-ir-mode rich}: Two-column layout with an IR source panel and a moving spotlight cursor alongside the stack view.
        \item \texttt{-{}-ir-mode rich-ssa}: Three-column layout (IR Source $\mid$ SSA Values $\mid$ Stack) showing computed results for binary operations, comparisons, and loads.
    \end{itemize}

    \item \textbf{CFG Traversal:} \texttt{-{}-cfg-animate} renders an animated walk through the control-flow graph using \texttt{CFGRenderer} (another \texttt{CommandDrivenScene} subclass) with DOT-derived node positions (\texttt{-{}-dot-cfg}) and a runtime trace (\texttt{-{}-import-trace}). If no trace is provided, a static DFS trace is derived automatically.

    \item \textbf{Export:} Static artifacts complement the animated output.
    \begin{itemize}
        \item \texttt{-{}-json}: Writes a~machine-readable scene graph (\texttt{scene\_graph.json}).
        \item \texttt{-{}-draw}: Writes \texttt{cfg\_main.dot} and, when Graphviz is available, a~rendered \texttt{cfg\_main.png}.
    \end{itemize}

    \item \textbf{Rendering:} Manim CE renders MP4 or GIF output. GIF conversion uses an \texttt{ffmpeg} palette workflow to reduce peak memory usage. The \texttt{-{}-preview} flag opens the rendered animation in a viewer.
\end{itemize}
